\documentclass[11pt]{article}
\usepackage[a4paper,margin=1in]{geometry}
\usepackage{microtype}
\usepackage{mathtools,amssymb,amsthm}
\usepackage{tikz-cd}
\usepackage[hidelinks]{hyperref}

\title{DaTra: Data Transformation Semantics As a Symmetrical Monoidal
Category on a Paginated Topos}
\author{\texttt{@qualiascript} \and ChatGPT}
\date{}

\theoremstyle{definition}
\newtheorem{definition}{Definition}[section]
\theoremstyle{plain}
\newtheorem{lemma}[definition]{Lemma}

\newcommand{\Set}{\mathsf{Set}}
\newcommand{\id}{\operatorname{Id}}
\newcommand{\Hom}{\operatorname{Hom}}
\newcommand{\PSh}{\operatorname{PSh}}

\begin{document}
\maketitle

\begin{abstract}
This file formalizes \texttt{datra.md}. The declarations are kept in the same
numbered order as the source document. The original prose is repeated
verbatim before the corresponding Lean declarations.

The implementation uses a chosen injection into $\mathbb{N}$ as the rank of an
element. This choice is isolated in \texttt{Dominion.rank}, so a later
replacement by accessible ordinals does not affect the categorical interfaces.
\end{abstract}

\section{Dominions}

\begin{definition}[The Category of Dominions]
The \textbf{Category of Dominions}, denoted $\mathsf{Dom}$, has as its objects
sets whose cardinality is at most $\omega_0$ (i.e. the smallest infinite
ordinal), that is, its objects are sets that have an injection to $\omega_0$,
and as its morphisms set-theoretic functions (i.e. morphisms in $\Set$).
\end{definition}

\begin{definition}[The Domanial Embedding Functor]
The \textbf{Domanial Embedding Functor}, denoted
$\mathsf{Emb} : \mathsf{Dom} \to \Set$, sends each dominion to its
corresponding set.
\end{definition}

\section{Domanial Insertions}

\begin{definition}[The Category of Domanial Insertions]
The \textbf{Category of Domanial Insertions}, denoted $\mathsf{DomIns}$, is
the category whose objects are dominions, that is, objects of $\mathsf{Dom}$,
and whose morphisms are monomorphisms in $\mathsf{Dom}$.
\end{definition}

\begin{definition}[The Category of Codomanial Insertions]
The \textbf{Category of Codomanial Insertions}, denoted $\mathsf{CoDomIns}$,
is the opposite category of $\mathsf{DomIns}$, that is,
\[
  \mathsf{CoDomIns} = \mathsf{DomIns}^{\mathrm{op}}.
\]
\end{definition}

\section{Domanial Consolidations}

\begin{definition}[The Category of Domanial Consolidations]
The \textbf{Category of Domanial Consolidations}, denoted $\mathsf{DomCon}$,
is the category whose objects are dominions and morphisms are order-preserving
epimorphisms in $\mathsf{Dom}$. More specifically, a morphism
$F : A \to B$ in $\mathsf{DomCon}$ has the property that for any $x,y : A$,
if $\lvert x\rvert < \lvert y\rvert$, then
$\lvert F(x)\rvert \leq \lvert F(y)\rvert$.
\end{definition}

\begin{definition}[The Category of Codomanial Consolidations]
The \textbf{Category of Codomanial Consolidations}, denoted
$\mathsf{CoDomCon}$, is the opposite category of $\mathsf{DomCon}$, that is,
\[
  \mathsf{CoDomCon} = \mathsf{DomCon}^{\mathrm{op}}.
\]
\end{definition}

\section{Chains}

\begin{definition}[Chain]
A \textbf{Chain} is a totally ordered thin category, taken skeletally.
\end{definition}

\begin{definition}[Short Chain]
A \textbf{Short Chain} is a chain that is finite, taken skeletally.
\end{definition}

\section{Transportation and Folios}

\begin{definition}[The Transportation Functor]
The \textbf{Transportation Functor}, denoted $\mathsf{Tra}$, is of type
$\mathsf{Tra} : \mathsf{DomCon} \to \Set$, and sends each morphism in
$\mathsf{DomCon}$ to its underlying set-theoretic surjection.
\end{definition}

\begin{definition}[Folio]
A \textbf{Folio} is a functor
$F : C \to (1 \downarrow \mathsf{CoDomCon})$, where
$1 \downarrow \mathsf{CoDomCon}$ is the coslice category, that preserves the
initial object.
\end{definition}

\section{Paginations}

\begin{definition}[The Category of Paginations]
The \textbf{Category of Paginations}, denoted $\mathsf{Pag}$, is the category
whose objects are pairs $(H,E)$, where
$H = \mathsf{Tra} \circ F^{\mathrm{op}}$,
$F : C \to \mathsf{CoDomCon}$ is a folio,
and $E$ is the category of elements of functor $H$. For any
$W : \mathsf{Pag}$, we denote $W_H$ as the first inclusion and $W_E$ as the
second inclusion. Morphisms in $\mathsf{Pag}$, for $F : X \to Y$ and
$X,Y : \mathsf{Pag}$, are given by functors $T : X_E \to Y_E$, so that the
following diagram commutes:
\[
\begin{tikzcd}[column sep=large,row sep=large]
x=(m,i) \arrow[r,"f"] \arrow[d,"T"']
  & \bigl(n,(X_H(f))(i)\bigr) \arrow[d,"T"] \\
(m',i') \arrow[r,"T(f)"']
  & \bigl(n',(Y_H(T(f)))(i')\bigr).
\end{tikzcd}
\]
\end{definition}

\section{Data Transformations}

\begin{definition}[The Category of Data Transformations]
The \textbf{Category of Data Transformations}, denoted $\mathsf{DaTra}$,
which we will also refer to as the \textbf{Category of Data Transformation
Sets}, or simply \textbf{DaTra Sets}, is the category whose objects are pairs
$(P,G)$, where $P : \mathsf{Pag}$ and $G$ is a functor
$G : P_E \to \mathsf{DomIns}$, or alternatively a presheaf
$G : P_E^{\mathrm{op}} \to \mathsf{CoDomIns}$. For
$X : \mathsf{DaTra}$, we denote
$X_P$ the first inclusion, $X_G$ the second inclusion, along with
$X_H = X_{P_H}$ and $X_E = X_{P_E}$. Morphisms in $\mathsf{DaTra}$, for
$T : X \to Y$, are pairs $(T_P,T_A)$, where $T_P : X_P \to Y_P$ is a
morphism in $\mathsf{Pag}$ and
$T_A : X_G \Rightarrow Y_G \circ T_E$ is a natural transformation, that is,
the following diagram commutes for all morphisms $f : x \to y$ in $X_P$:
\[
\begin{tikzcd}[column sep=huge,row sep=large]
X_G(x) \arrow[r,"X_G(f)"] \arrow[d,"(T_A)_x"']
  & X_G(y) \arrow[d,"(T_A)_y"] \\
Y_G(T_E(x)) \arrow[r,"Y_G(T_E(f))"']
  & Y_G(T_E(y)).
\end{tikzcd}
\]
\end{definition}

\begin{lemma}[Transformation Lemma]
The \textbf{Transformation Lemma} states that $\mathsf{DaTra}$ is a topos.

\begin{proof}[Proof sketch]
Let $\PSh(B) = B^{\mathrm{op}} \to \Set$ be the fixed category of underlying
presheaves. The forgetful functor
$U : \mathsf{DaTra} \to \PSh(B)$ has a \textbf{Minimal Pagination Functor}
$\mathsf{MinPag} : \PSh(B) \to \mathsf{DaTra}$. The adjunction
$\mathsf{MinPag} \dashv U$ has invertible unit and counit, so
$\mathsf{DaTra} \simeq \PSh(B)$, and $\mathsf{DaTra}$ is a topos.
\end{proof}
\end{lemma}

\begin{definition}[Cardinality of a DaTra Set]
The \textbf{Cardinality of a DaTra Set} $D : \mathsf{DaTra}$ is the
cardinality of the origin of $D_H : C \to \Set$, that is, the number of
objects of the small chain $C$. It is denoted as $\lvert D\rvert$.
\end{definition}

\begin{definition}[Extent of a DaTra Set]
The \textbf{Extent of a DaTra Set} $D : \mathsf{DaTra}$ is the value at
$D_G(0,0)$, that exists as folios preserve initial objects. It is denoted as
$\mathsf{Ex}(D)$. Since the objects of $\mathsf{DomIns}$ are dominions,
$\mathsf{Ex}(D) : \mathsf{Dom}$ for any $D : \mathsf{DaTra}$.
\end{definition}

\begin{definition}[Territory of a DaTra Set]
The \textbf{Territory of a DaTra Set} $D : \mathsf{DaTra}$ is the indexed
collection of dominions $\mathsf{Ter}(D) : M \to \mathsf{Dom}$, where
$M = D_G(\lvert D\rvert - 1)$, so that
$\mathsf{Ter}(D)(k) = D_G(\lvert D\rvert - 1,k)$.
\end{definition}

\begin{definition}[$n$th Region of a DaTra Set]
The \textbf{$n$th Region of a DaTra Set} $D : \mathsf{DaTra}$ is the
dominion $\mathsf{Ter}(D)(n)$ for $n : M$, noting that the indexing starts at
$0$.
\end{definition}

\section{Transposals and Traversals}

\begin{definition}[The Category of Data Transposals]
The \textbf{Category of Data Transposals}, denoted $\mathsf{DaTrap}$, is the
wide subcategory of $\mathsf{DaTra}$ whose morphisms $F : X \to Y$ have the
property that $F_E$ is a faithful functor.
\end{definition}

\begin{definition}[The Category of Data Traversals]
The \textbf{Category of Data Traversals}, denoted $\mathsf{DaTrav}$, is the
wide subcategory of $\mathsf{DaTrap}$ whose morphisms $F : X \to Y$ have the
following properties: for any $x=(k,i)$ and $y=(k,j)$ with $i<j$, let
$x'=F_E(x)=(m,i')$, $y'=F_E(y)=(n,j')$, and let
$k'=\min(m,n)$; then
\[
  Y_H(m \to k')(i') < Y_H(n \to k')(j').
\]
\end{definition}

\section{Maps and Cartography}

\begin{definition}[The Category of Data Transformation Maps]
The \textbf{Category of Data Transformation Maps}, denoted
$\mathsf{DaTraMap}$, is the full subcategory of $\mathsf{DaTra}$ that includes
all DaTra sets with the property that their extent is isomorphic to the
coproduct of its regions. That is, for $D : \mathsf{DaTra}$, let
$F = \mathsf{Emb} \circ \mathsf{Ter}(D) : M \to \Set$, and let
$\mathsf{Sum}(D)$ be the coproduct on diagram $F$; then
$D : \mathsf{DaTraMap}$ if and only if
$\mathsf{Sum}(D) \simeq \mathsf{Ex}(D)$.
\end{definition}

\begin{definition}[The DaTra Map Inclusion Functor]
The \textbf{DaTra Map Inclusion Functor},
$\mathsf{DaTraMapInc} : \mathsf{DaTraMap} \to \mathsf{DaTra}$, sends each
DaTra map to its equivalent DaTra set.
\end{definition}

\begin{definition}[The Category of Data Traversal Maps]
The \textbf{Category of Data Traversal Maps}, or simply \textbf{DaTrav Maps},
denoted $\mathsf{DaTravMap}$, is the wide subcategory of
$\mathsf{DaTraMap}$ so that a morphism $f : x \to y$ of $\mathsf{DaTraMap}$
is a morphism of $\mathsf{DaTravMap}$ if and only if
$\mathsf{DaTraMapInc}(f)$ is a morphism of $\mathsf{DaTrav}$.
\end{definition}

\begin{definition}[The DaTrav Map Inclusion Functor]
The \textbf{DaTrav Map Inclusion Functor},
$\mathsf{DaTravMapInc} : \mathsf{DaTravMap} \to \mathsf{DaTrav}$, is the
canonical restriction of $\mathsf{DaTraMapInc}$ on the morphisms of
$\mathsf{DaTravMap}$.
\end{definition}

\begin{definition}[The Charting Functor]
The \textbf{Charting Functor}, if it exists, is defined as the right adjoint
to $\mathsf{DaTravMapInc}$, and it is denoted
$\mathsf{Chr} : \mathsf{DaTrav} \to \mathsf{DaTravMap}$.
\end{definition}

\begin{lemma}[Cartography Lemma]
The \textbf{Cartography Lemma} states that $\mathsf{Chr}$ exists.

\begin{proof}[Proof sketch]
$\mathsf{Chr}$'s action on objects $X : \mathsf{DaTrav}$ is sending it to the
universal arrow from $\mathsf{DaTravMapInc}$ to $X$, so that we denote
$X' = \mathsf{DaTravMapInc}(\mathsf{Chr}(X))$, along with $H : X' \to X$,
with $X'$ terminal among objects with this property. The solution is given by
the DaTra map so that $\mathsf{Ter}(X') \simeq \mathsf{Ter}(X)$, which is
unique as $\mathsf{DaTrav}$ preserves orders, and terminal as $H_E$ is
faithful.
\end{proof}
\end{lemma}

\section{Coalitions}

\begin{definition}[The Coalizing Functor]
The \textbf{Coalizing Functor}, denoted
$\mathsf{Coa} : \mathsf{DaTrav} \to \mathsf{Dom}$, is defined as
$\mathsf{Coa} = \mathsf{Ex} \circ \mathsf{Chr}$.
\end{definition}

\begin{definition}[The Domanial Inclusion Functor]
The \textbf{Domanial Inclusion Functor}, denoted
$\mathsf{DomInc} : \mathsf{Dom} \to \mathsf{DaTrav}$, is the functor that is
left adjoint to $\mathsf{Coa}$. That is, for every $X : \mathsf{Dom}$ and
$Y : \mathsf{DaTrav}$,
\[
  \Hom_{\mathsf{DaTrav}}\bigl(\mathsf{DomInc}(X),Y\bigr)
  \simeq
  \Hom_{\mathsf{Dom}}\bigl(X,\mathsf{Coa}(Y)\bigr),
\]
naturally in $X$ and $Y$, so that
$\mathsf{Coa}(\mathsf{DomInc}(X))=X$.
\end{definition}

\begin{definition}[The Coalition of a DaTra Object]
The \textbf{Coalition of a DaTra Object}, given an object $X$ of
$\mathsf{DaTra}$, is $\mathsf{Coa}(X')$, where $X' : \mathsf{DaTrav}$ is the
same object viewed as an object of $\mathsf{DaTrav}$.
\end{definition}

\begin{definition}[The Empty Map]
The \textbf{Empty Map} is the DaTra set that $\mathsf{DomInc}$ sends the
initial object of $\mathsf{Dom}$, $0 : \mathsf{Dom}$, to. It is denoted as
$I=\mathsf{DomInc}(0)$, and it is initial in $\mathsf{DaTra}$.
\end{definition}

\section{Horizontal Sums}

\begin{definition}[The Horizontal Sum Bifunctor]
The \textbf{Horizontal Sum Bifunctor}, if it exists, is denoted as
$\mathsf{HorSum} : \mathsf{DaTrav} \times \mathsf{DaTrav}
\to \mathsf{DaTrav}$, or using infix notation using the $+_{<}$ symbol as
$\mathsf{DaTrav} +_{<} \mathsf{DaTrav} \to \mathsf{DaTrav}$, and is defined
as follows: given two morphisms in $\mathsf{DaTrav}$, $F : X \to Y$ and
$F' : X \to Y$, let $H : G \times G' \to F+F'$ be the universal arrow from
$\mathsf{DaTrav} \times \mathsf{DaTrav}$ to $F+F'$ so that for projections
$H_1 : G \to F+F'$ and $H_2 : G' \to F+F'$, for
$(m,i)=H_{1_E}(0,0)$ and $(m',i')=H_{2_E}(0,0)$, we have $m<2$, $m'<2$,
$i=0$, and furthermore, if $m=1$, then $i'=1$. Then $F+_{<}F'=H$. It
remains to be shown $H$ exists and is unique for all morphisms $F,F'$ of
$\mathsf{DaTrav}$.
\end{definition}

\begin{lemma}[Horizontal Lemma]
The \textbf{Horizontal Lemma} states that $\mathsf{HorSum}$ exists.

\begin{proof}[Proof sketch]
For $F+_{<}F'$, if $F'=I$, then $H : I \to F$, and $I$ is initial so $H$
is unique, so that $H \simeq H_1 \simeq H_2$ and
$(m,i)=H_E(0,0)=(0,0)$, so that $i=0$ and as $m=0$, the condition on $i'$
need not be checked. The case for $F=I$ is dual. If both $F$ and $F'$ are
distinct from $I$, for $(m,i)=H_{1_E}(0,0)$, if $m=0$, then $i=0$ and
$\mathsf{Ex}(F)=\mathsf{Ex}(F)+\mathsf{Ex}(F')$, but as $F'$ is not $I$,
this cannot hold, so that $m=1$, which implies $i'=1$, so that the solution
is given by $(m,i)=(1,0)$ and $(m',i')=(1,1)$.
\end{proof}
\end{lemma}

\section{Symmetric Monoidal Structure}

\begin{definition}[The Data Traversal Monoidal Category]
The \textbf{Data Traversal Monoidal Category}, denoted $\mathsf{DaTravMon}$,
if it exists, is the symmetric monoidal category with $\mathsf{DaTrav}$ as the
underlying category, $+_{<}$ as the tensor product and $I$ as the identity,
so that we denote
\[
  \mathsf{DaTravMon}=(\mathsf{DaTrav},+_{<},I).
\]
As for $X : \mathsf{DaTrav}$, $X=X+I=I+X$ is strict, the left and right
unitors are given by identity. It remains to be shown that there is a braider
functor whose double application yields identity, and an associator functor so
that the pentagon, triangle and hexagon identities commute.
\end{definition}

\begin{definition}[The Data Traversal Braider]
The \textbf{Data Traversal Braider}, denoted
$\mathsf{Brd}_{X_Y} : X+_{<}Y \to Y+_{<}X$ for
$X,Y : \mathsf{DaTrav}$, is defined as follows: if $X=I$ or $Y=I$,
$\mathsf{Brd}_{X_Y}=\id$. Otherwise, let $G : H \to X+_{<}Y$ be the
universal arrow from $\mathsf{DaTrap}$ to $X+_{<}Y$ so that
$G_E(1,0)=(1,1)$ and $G_E(1,1)=(1,0)$. Trivially, $H=Y+_{<}X$, so that we
let $\mathsf{Brd}_{X_Y}(X+_{<}Y)=H$, and trivially,
$\mathsf{Brd}_{Y_X}\mathsf{Brd}_{X_Y}=\id$.
\end{definition}

\begin{definition}[The Data Traversal Associator]
The \textbf{Data Traversal Associator}, denoted
$\mathsf{Asoc}_{X_{Y_Z}} : (X+_{<}Y)+_{<}Z \to X+_{<}(Y+_{<}Z)$, is
defined as follows: if $X=I$, $Y=I$ or $Z=I$, then
$\mathsf{Asoc}_{X_{Y_Z}}=\id$. Otherwise, denote
$R=(X+_{<}Y)+_{<}Z$ and let $G : H \to R$ be the universal arrow from
$\mathsf{DaTrap}$ to $R$ so that $G_E(2,0)=(1,0)$ and there exists a
morphism $G' : (Y+_{<}Z) \to H$ in $\mathsf{DaTrav}$ with
$G'_E(0,0)=(1,1)$. As $G_E$ is faithful and $G_A=\id$, $G$ is invertible in
$\mathsf{DaTrap}$, and by extension in $\mathsf{DaTra}$.
\end{definition}

\begin{lemma}[Serenity Lemma]
The \textbf{Serenity Lemma} states that $\mathsf{DaTravMon}$ exists.

\begin{proof}[Proof sketch]
As both $\mathsf{Asoc}$ and $\mathsf{Brd}$ send objects to isomorphic objects
in $\mathsf{DaTra}$, it is clear that the pentagon, triangle and hexagon
identities commute, thus $\mathsf{DaTravMon}$ is a symmetric monoidal
category.
\end{proof}
\end{lemma}

\end{document}
